\section{Introduction}\label{sec:introduction}

Arnold's fixed-point conjecture compares Hamiltonian dynamics with the
critical-point theory of smooth functions. For a sufficiently short
autonomous Hamiltonian flow on a closed symplectic manifold, its fixed points are
exactly the critical points of its Hamiltonian. The conjecture asks
whether an arbitrary Hamiltonian diffeomorphism must retain the minimum
number of critical points forced by the topology of the manifold.
Arnold's formulation explicitly includes the count of geometrically
distinct points~\cite[Section~2]{Arnold1986}. We construct a smooth
counterexample on the complex quadric threefold, which also violates the
weaker rational cup-length bound.

\paragraph{The two degenerate bounds.}
For a closed smooth manifold $M$, write
\[
 \CritSet(h)=\{q\in M:dh_q=0\},\qquad
 \Crit(M)=\min_{h\in C^\infty(M,\R)}\#\CritSet(h).
\]
The minimum here permits degenerate critical points, whereas the
\emph{Morse number} minimizes only over Morse functions. Let
$\operatorname{cuplength}(M;\mathbb Q)$ be one plus the maximum number
of positive-degree rational cohomology classes with nonzero cup product.
Thus our cup-length convention includes the unit. The
Lusternik--Schnirelmann inequality gives
\[
 \operatorname{cuplength}(M;\mathbb Q)\le\Crit(M);
\]
Section~\ref{sec:critical-bound} recalls the cup-product argument in a
form that applies to arbitrary smooth functions.

A Hamiltonian diffeomorphism of a symplectic manifold $(M,\omega)$ is
the time-one map of a smooth time-dependent Hamiltonian flow. We write
$\Ham(M,\omega)$ for these maps, use the convention
$\iota_{X_{H_t}}\omega=dH_t$, and set
$\Fix(\phi)=\{q\in M:\phi(q)=q\}$.
The \emph{critical-number Arnold bound} asserts that
\begin{equation}\label{eq:arnold-bound}
 \#\Fix(\phi)\ge\Crit(M)
\end{equation}
for every $\phi\in\Ham(M,\omega)$ on every closed connected symplectic
manifold. Its weaker \emph{rational cup-length} form replaces $\Crit(M)$
by $\operatorname{cuplength}(M;\mathbb Q)$. Neither formulation assumes
nondegeneracy; see Golovko~\cite[Conjectures~1.4 and~1.6]{Golovko2020}.
A fixed point $q$ is nondegenerate when $1$ is not an eigenvalue of
$d\phi_q$.

Let
\begin{equation}\label{eq:quadric}
 Q^3=\bigl\{[z_0:z_1:z_2:z_3:z_4]\in\CP^4:
                 z_0^2+z_1^2+z_2^2+z_3^2+z_4^2=0\bigr\},
\end{equation}
with the restricted Fubini--Study symplectic form, normalized in
Section~\ref{sec:quadric}.

\begin{theorem}\label{thm:counterexample}
There is a smooth Hamiltonian diffeomorphism $\phi$ of the closed connected
symplectic six-manifold $Q^3$ such that
\begin{equation}\label{eq:counterexample}
 \#\Fix(\phi)=3<4=\Crit(Q^3)
                  =\operatorname{cuplength}(Q^3;\mathbb Q).
\end{equation}
At least one fixed point of $\phi$ is degenerate.
\end{theorem}

This disproves both unrestricted smooth bounds. The count includes all
fixed points, so imposing a contractibility condition on their
Hamiltonian trajectories cannot restore either inequality. The
nondegenerate homological Arnold inequalities concern a different
hypothesis and remain compatible with this example.

The count of three is sharp. Gong~\cite[Section~1.1.1]{Gong2025}
obtains at least $n$ fixed points for every Hamiltonian diffeomorphism
of the standard complex quadric of complex dimension $n\ge2$.
His normalization of the Fubini--Study form differs from ours by a
positive scalar, which does not change the Hamiltonian group: rescaling
the form and a generating Hamiltonian by the same scalar preserves its
vector field. Thus our example attains his bound for $n=3$.

\paragraph{Historical context and related work.}
The positive cases explain why the unrestricted comparison is natural.
Conley and Zehnder proved that every Hamiltonian diffeomorphism of the
standard symplectic torus $T^{2n}$ has at least $2n+1$ fixed
points~\cite{ConleyZehnder1983}, and Fortune obtained the bound $n+1$
on standard $\CP^n$~\cite{Fortune1985}. Hofer obtained cup-length bounds without
nondegeneracy when $\pi_2(M)=0$~\cite[Theorem~4]{Hofer1988}.
Rudyak and Oprea proved the full critical-number bound when both the
symplectic class and the first Chern class vanish on $\pi_2(M)$
\cite[Corollary~4.3]{RudyakOprea1999}. The quadric has spheres of
positive symplectic area and lies outside these latter hypotheses.

Floer's pseudoholomorphic-curve methods established homological Morse
inequalities for nondegenerate Hamiltonian fixed points, including the
monotone setting~\cite{Floer1989}. The lower bound by the sum of rational
Betti numbers was subsequently established for nondegenerate Hamiltonian
diffeomorphisms on every closed symplectic manifold by Fukaya and Ono,
Liu and Tian, and Ruan~\cite{FukayaOno1999,LiuTian1998,Ruan1999}.
The distinction between these bounds
and the degenerate formulations above is essential; see
\cite[Theorem~1.3 and Conjectures~1.4 and~1.6]{Golovko2020}.
Ma~\cite[Theorem~1.5]{Ma2013} asserted that the minimum number of fixed
points over all Hamiltonian diffeomorphisms always equals $\Crit(M)$.
Theorem~\ref{thm:counterexample} contradicts that assertion.

Regularity also matters. Buhovsky, Humili\`ere, and Seyfaddini constructed
Hamiltonian \emph{homeomorphisms} with a single fixed point on every closed
connected symplectic manifold of real dimension at least four
\cite{BHS2018}. Our construction stays in the smooth Hamiltonian group.
In a preprint posted September 27, 2026, Jiao~\cite[Theorem~1.1]{Jiao2026}
reports a smooth Hamiltonian diffeomorphism of the equal-area product
$S^2\times S^2$ with two fixed points, using local saddle merging and
mixed generating functions. The construction here instead uses a
finite-order map on a six-dimensional quadric. None of these
fixed-point results is an input to our construction.

\paragraph{Proof mechanism.}
The key is to transfer a critical-point count from a fixed submanifold
to the fixed-point count of a nearby Hamiltonian map. Let $A$ be a
finite-order Hamiltonian map with fixed submanifold $F$, whose fixed
tangent vectors are exactly those tangent to $F$. Extend a function
$f$ on $F$ to an $A$-invariant function $K$ on the ambient manifold,
and let $B_s$ be its Hamiltonian flow. If $A^m=\id$, invariance gives
\[
 (A\circ B_\varepsilon)^m=B_{m\varepsilon}.
\]
For sufficiently small $\varepsilon>0$, the right side has only
stationary fixed points. Any fixed point of $A\circ B_\varepsilon$
is therefore a critical point of $K$ fixed by $A$. Finite-group
averaging makes ambient criticality on $F$ equivalent to criticality
of $f$, giving
\[
 \Fix(A\circ B_\varepsilon)=\CritSet(f).
\]
This criticality step is the finite-group case of Palais's principle
of symmetric criticality~\cite[Section~5]{Palais1979}; the power identity
excludes additional fixed points away from $F$.

For the quadric, take the Hamiltonian involution
\[
 A[z_0:z_1:z_2:z_3:z_4]=[z_0:-z_1:-z_2:-z_3:-z_4].
\]
Its fixed locus is the quadric surface $Q^2\cong S^2\times S^2$.
This product admits functions with only three critical points, as
recorded by Mare~\cite[Proposition~3.2.5]{Mare1997} through an extension
of Takens's sphere-product construction~\cite{Takens1968}. We give an
explicit formula and compute its entire critical set. The isolated
degenerate critical point becomes a degenerate Hamiltonian fixed point.
In contrast, symplectic volume forces every smooth function on the
ambient $Q^3$ to have at least four critical points, and a Hamiltonian
rotation realizes that minimum.

Section~\ref{sec:reduction} proves the finite-order transfer, including
the short-period estimate and invariant extension.
Section~\ref{sec:quadric} constructs the Hamiltonian involution and
its tangent splitting. Section~\ref{sec:three-points} gives the
three-critical-point function. Section~\ref{sec:critical-bound}
proves the critical-number and rational cup-length equalities on $Q^3$,
and Section~\ref{sec:conclusion} assembles the construction and
verifies degeneracy.
